\chapter{Conclusion and Future Work}

\section{Conclusion}

This project set out to build a remote desktop system that provides the
reachability users expect from commercial tools without asking them to trust an
opaque relay, and that additionally supports genuine unattended operation. Both
aims were met.

The system pairs a lightweight Rust signaling server with direct, end-to-end
encrypted WebRTC media. The server authenticates every endpoint with an Ed25519
challenge--response over a closed, code-gated enrollment scheme, yet is
structurally unable to decrypt the sessions it arranges because the DTLS-SRTP
keys are negotiated only between the peers. This separation --- the server knows
\emph{who} is connecting but never \emph{what} is sent --- is the core
contribution and the property that distinguishes \updesk{} from relay-mediated
products.

The second contribution is a headless native host. By capturing the framebuffer
directly rather than through a browser API, it operates with no window, no
gesture and no picker; by encoding to H.264 and publishing on a native WebRTC
track it reuses the same secure transport as the attended host; and by
registering as a Windows service it is present from boot and can reach the active
desktop. Native input injection makes it a full control host and not merely a
viewer.

Evaluation on both a local network and the public Internet confirmed direct
peer-to-peer connectivity across the tested NAT configurations, interactive
latency suitable for real use, and unattended service operation that survived
reboot and session switching. A real defect in DTLS curve negotiation against
modern browsers was diagnosed and repaired in a vendored patch --- a repair that
was only possible because every layer of the stack was open to inspection.

\section{Future Work}

Several directions would extend the system.

\begin{itemize}
  \item \textbf{Complete secure-desktop capture.} The SYSTEM service provides the
        foundation for capturing the Winlogon secure desktop; dedicated
        implementation and testing of that path would let a controller drive UAC
        prompts and the login screen, closing the last gap with RDP.
  \item \textbf{Hardware-accelerated, low-latency encoding.} Replacing the
        software H.264 path with a hardware encoder configured for low latency
        would reduce wide-area delay and free CPU on the host.
  \item \textbf{Audio and clipboard.} Forwarding host audio and synchronising the
        clipboard are natural additions over the existing data channel, deferred
        deliberately to keep the video and control paths correct first.
  \item \textbf{File transfer.} A negotiated file-transfer channel would round out
        the support use case.
  \item \textbf{Multi-monitor support.} Extending capture beyond the primary
        display, with monitor selection at the controller.
  \item \textbf{Cross-platform hosts.} The native capture and input layers are the
        platform-specific parts; Linux and macOS hosts would follow the same
        architecture with different capture and injection back ends.
  \item \textbf{Mature mobile clients.} The repository contains exploratory
        Android host and controller work that could be developed into supported
        clients.
\end{itemize}

Taken together, these would turn a system that already demonstrates the central
security and unattended-operation ideas into a complete alternative to the
proprietary tools it was conceived against.
